\section{Theorie}
\label{sec:Theorie}
\subsection{Fourier-Theorie}
Das Fouriersche Theorem besagt, dass sich jedes periodische
Signal $f(t)$ als eine Reihe von Sinus- und Cosinusfunktionen der Form
\begin{equation}
    \label{eqn:Fourier-Reihe}
    f(t)=\frac{1}{2}a_0+\sum_{n=1}^{\infty}\left(a_n \cos(n\frac{2\symup{\pi}}{T}t)+b_n \sin(n\frac{2\symup{\pi}}{T}t)\right)
\end{equation}
beschreiben lässt. Wobei T die Periodendauer des Signals ist.
Mit der Periodendauer T lässt sich zur Grundfrequenz $\frac{1}{\symup{T}}$ eine Kreisfrequenz $\omega=\frac{2\symup{\pi}}{T}$
definieren. Die Koeffizienten $a_n$ und $b_n$ beschreiben die Anteile aller
Oberwellen, Schwingungen mit einem ganzzahligen Vielfachen der Grundfrequenz,
die im Signal enthalten sind. Analytisch lassen Sie sich mit den Formeln
\begin{align}
    a_n & =\frac{2}{T}\int_0^Tf(t)\cos{(n\frac{2\symup{\pi}}{T}t)}\symup{d}t &
    b_n & =\frac{2}{T}\int_0^{T}f(t)\sin{(n\frac{2\symup{\pi}}{T}t)}\symup{d}t
    \label{eqn:Koeffizienten}
\end{align}
berechnen. Ein Anteil für n=0 entspricht dabei einem sogenannten Offset.
Das Signal ist im Mittel nicht 0, sondern schwingt um einen von 0
versetzten Wert.\\

\noindent Diese Reihendarstellung des Signals konvergiert gleichmäßig
gegen das zu beschreibende Signal $f(t)$, wenn $f(t)$ stetig ist. An Stellen der Unstetigkeit
tritt das Gibbsche Phänomen auf. Demnach kommt es im Bereich der Unsteigkeit zu einer
Überschwingung, die keinem experimentellen Fehlern zugrunde liegt.
\cite{V351}
\subsection{Fourier-Analyse}
Bei der Fourier-Analyse wird ein zeitlich periodisches Signal auf
seine Frequenzanteile geprüft. Dazu werden analoge und digitale Verfahren entworfen.
Ein analoges Verfahren basiert auf den Eigenschaften eines elektrischen Schwingkreises. Bei einer speziellen Frequenz,
der Resonanzfrequenz, besitzt dieser einen merklichen Überschwinger. Wird ein Signal an den Eingang einer solchen Schaltung
gelegt und die Grundfrequenz auf Resonanzfrequenz gestellt, können die Amplitudenverhältnisse der einzelnen Oberwellen zur Grundfrequenz gemessen werden.
Da bei neu eingestellter Frequenz $\nu/n$, die n-te Oberwelle Resonanz am Ausgang hervorruft.\\

\noindent Bei digitalen Verfahren wird das Signal quantisiert. Mit der sogenannten
Abtastfrequenz, die mindestens doppelt so groß wie die höchste zu messende Oberwellenfrequenz sein soll,
werden Messwerte genommen und dann mit einer Fast-Fourier-Transformation (FFT) in den
Frequenzbereich übertragen, wo die Frequenzanteile dann abgelesen werden können.\\

\noindent Die Amplitudenverhältnisse werden teilweise in Dezibel angegeben. Bei gegebener Referenzspannung
$\symup{U_0}$ (z.B. Spannung der Grundfrequenz) ist das Verhältnis in Dezibel [ dB ] gegeben durch
\begin{equation}
    \left|\frac{U_n}{U_0}\right|_{\symup{dB}}=20\log{\frac{U_n}{U_0}},
    \label{eqn:Dezibel}
\end{equation}
wobei mit dem 10er-Logarithmus gearbeitet wird. Das richtige Spannungsverhältnis aus einer dB Angabe
kann dann nach Umstellen mit
\begin{equation}
    \frac{U_n}{U_0}=10^{\frac{1}{20}\cdot\left|\frac{U_n}{U_0}\right|_{\symup{dB}}}
\end{equation}
berechnet werden \cite{V351}, \cite{NeuesGP}.
\subsection{Fourier-Synthese}
In der Fouriersynthese werden periodische Signale aus bereits bestehenden Oberwellen
zusammengesetzt. Das heißt die Amplituden, also die Fourierkoeffizienten,
lassen sich so einstellen, dass ein gewünschtes Ausgangssignal entsteht (z.B. Rechtecksignal).
\cite{V351}
\subsection{Vorbereitung: Fourierreihe zu 3 periodischen Signalen}
\subsubsection{Rechtecksignal}
Ein Rechtecksignal der Art
\begin{equation}
    f_0(t)=
    \begin{cases}
        h  & 0\leq t < \frac{\symup{L}}{2}  \\
        -h & \frac{\symup{L}}{2} \leq t <\symup{L} \\
        0  & \text{sonst}           \\
    \end{cases}
    \label{eqn:Rechteck}
\end{equation}
wird periodisch fortgesetzt, also um nL verschoben unendlich häufig addiert.
Es ist sofort einzusehen, dass es keinen Offset gibt, da es gleiche Anteile über und
unter Null besitzt. Außerdem ist das Signal eine ungerade Funktion um die Y-Achse,
also gehen nur Sinusterme mit in die Rechnung ein. Für die Koeffizienten $b_n$
ergibt sich:
\begin{equation}
    b_n=\frac{2}{\symup{L}}\left(\int_0^{\frac{\symup{L}}{2}}h\sin{\frac{2\symup{\pi}nt}{\symup{L}}}\symup{d}t+\int_{\frac{\symup{L}}{2}}^{\symup{L}}(-h)\sin{\frac{2\symup{\pi}nt}{\symup{L}}}\symup{d}t\right)
    =\frac{2h}{\pi n}\left(1-(-1)^n\right)
    =\frac{4h}{\pi(2n-1)}  , n\in \symbb{N}
    \label{eqn:Rechteck-Koeff}
\end{equation}
\subsubsection{Sägezahn}
Eine Sägezahnfunktion mit fallender Flanke kann mit der Funktion
$f(t)=h(1-\frac{2t}{\symup{L}}), 0\leq t<L$ beschrieben werden. Sie hat ebenfalls
keinen Offset, da sie linear abfällt und die Nullstelle genau in der Intervallmitte liegt. Die Funktion
ist ebenfalls ungerade, weshalb nur Sinusterme auftauchen werden.
Die Koeffizienten $b_n$ ergeben sich dann mit partieller Integration zu:
\begin{equation}
    b_n=\frac{2}{\symup{L}}\left(\int_0^{\symup{L}}h\left(1-\frac{2t}{\symup{L}}\right)\sin{\frac{2\symup{\pi}nt}{\symup{L}}}\symup{d}t\right)
    =\frac{h}{\pi n}\left(1-1+2-\frac{2h}{L}\left[\frac{\sin{\frac{2\symup{\pi}nt}{\symup{L}}}}{\frac{2\symup{\pi}n}{\symup{L}}}\right]_0^{\symup{L}}\right)
    =\frac{2h}{\pi n}  , n\in \symbb{N}
    \label{eqn:Sägezahn-Koeff}
\end{equation}
\subsubsection{Dreieck}
Ein Dreieckssignal der Art
\begin{equation}
    f_0(t)=
    \begin{cases}
        (t+\frac{\symup{L}}{2})\frac{2}{\symup{L}}h & \frac{-\symup{L}}{2}\leq t < 0 \\
        (\frac{\symup{L}}{2}-t)\frac{2}{\symup{L}}h & 0 \leq t < \frac{\symup{L}}{2} \\
        0                                                   & \text{sonst}                   \\
    \end{cases}
    \label{eqn:Dreieck}
\end{equation}
wird periodisch fortgesetzt. Das Signal besitzt einen Offset, da das Signal dauerhaft
größer 0 ist. Außerdem ist das Signal gerade um 0. Daher gehen nur Cosinusfunktionen
in die Reihe mit ein. Im Folgenden wird der Offset ignoriert. Es werden nur die schwingenden Anteile
berücksichtigt, da nur diese für die Dreieckform relevant sind. Die Koeffizienten ergeben
sich mit partieller Integration und der Symmetrie zur Intervallmitte zu
\begin{equation}
    \begin{split}
        b_n&=\frac{2}{\symup{L}}\int_{\frac{-\symup{L}}{2}}^{\frac{\symup{L}}{2}}f_0(t)\cos{\frac{2\symup{\pi}nt}{\symup{L}}}\symup{d}t
        =\frac{4}{\symup{L}}\int_0^{\frac{\symup{L}}{2}}(\frac{\symup{L}}{2}-t)\frac{2}{\symup{L}}h\cos{\frac{2\symup{\pi}nt}{\symup{L}}}\symup{d}t\\
        &=\frac{2h}{(\pi n)^2}\left(1-(-1)^n\right)
        =\frac{4h}{(\pi(2n-1))^2}  , n\in \symbb{N}
    \end{split}
    \label{eqn:Dreieck-Koeff}
\end{equation}
(\textit{Die vollständige Berechnung befindet sich im Anhang})


